
\[
K_\bullet(X')\xrightarrow{i_*}K_\bullet(X)\xrightarrow{j^*}K_\bullet(U)\longrightarrow 0
\]
\uline{est
exacte}.

Si $\underline{F}'$ est un $\underline{O}_{X'}$-Module cohérent, $i_*(\underline{F}')$ est un $\underline{O}_X$-Module dont le support est contenu dans $X'$; par conséquent sa restriction à $U$ est nulle, et $j^*\circ i_*(\cl(\underline{F}'))=0$, d'où, les éléments de la forme $\cl(\underline{F}')$ engendrant $K_\bullet(X')$, $j^*\circ i_*=0$. On en déduit un homomorphisme
\[
 \rho:K_\bullet(X)/\operatorname{Im}K_\bullet(X')\longrightarrow K_\bullet(U).
\]
On va définir un homomorphisme inverse de $\rho$. Soit $\underline{G}$ un $\underline{O}_U$-Module cohérent. Comme $X$ est noethérien, il existe un $\underline{O}_X$-Module cohérent $\underline{F}$ dont la restriction à $U$ est isomorphe à $\underline{G}$. Si $\underline{F}'$ est un second $\underline{O}_X$-Module cohérent dont la restriction à $U$ est isomorphe à $\underline{G}$, alors on a dans $K_\bullet(X)$ :
\[
 \cl(\underline{F})-\cl(\underline{F}')\in \operatorname{Im}K_\bullet(X').
\]
Pour le voir, on considère $\underline{G}$ comme un sous-faisceau de $j^*(\underline{F}\oplus \underline{F}')$ au moyen de l'application diagonale; $\underline{G}$ se prolonge alors en un sous-faisceau $\underline{F}''$ de $\underline{F}\oplus \underline{F}'$, et on est ramené à montrer que
\[
 \cl(\underline{F})-\cl(\underline{F}'')\in \operatorname{Im}K_\bullet(X'),
\]
et l'analogue pour $\underline{F}'$. Soient alors $\underline{H}$ et $\underline{K}$ le noyau et le conoyau de l'homomorphisme composé :
\[
\underline{F}''\longrightarrow \underline{F}\oplus \underline{F}'\longrightarrow \underline{F} ;
\]
cet homomorphisme se réduit à l'identité de $\underline{G}$ sur $U$, et par suite $\underline{H}$ et $\underline{K}$ sont à support contenu dans $X'$. Si $\underline{I}$ est l'idéal de $X'$ dans $X$, il en résulte qu'il existe des entiers $m$ et $n$ tels que $\underline{I}^m\cdot \underline{H}=0$, $\underline{I}^n\cdot \underline{K}=0$ (EGA I 9.3.4). On peut donc dévisser $\underline{H}$ et $\underline{K}$ en extensions de Modules
